\section*{Общие сведения}
\addcontentsline{toc}{section}{Общие сведения}
В работе последовательно реализованы сбор корпуса, поисковый робот,
токенизация, морфологическая нормализация, оценка качества, построение
обратного индекса, булев и цитатный поиск, сжатие, прыжки по спискам
постингов и ранжирование TF-IDF. Корпус содержит документацию и обсуждения
по программированию, операционным системам, компиляторам и системным
библиотекам. Основные компоненты написаны на C++20, обвязка и робот на Python.

Нумерация соответствует списку лабораторных курса «Информационный поиск»
\cite{course}. Выполнены работы №1--8, 10--13 и 17.
Каждый раздел содержит задание, описание метода, журнал выполнения,
план тестирования, результаты и выводы.

Корпус содержит 1\,679\,971 документ и 17\,236\,388\,223 байта текста.
Для опыта с законом Ципфа отдельно выделены первые 10\,000 документов.
Это техническая выборка, а не случайная репрезентативная выборка всего корпуса.
Размеры, статистика токенизации и индексации относятся к указанному
набору данных; сравнение сжатия и прыжков выполняется при одинаковой
нормализации документов и запросов.

Качество поиска оценено на 30 запросах для четырёх систем.
Размечены все 426 уникальных пар «запрос, документ» объединённого пула
первых пяти результатов. Оценки выставлены ассистентом по содержимому
документов и информационным потребностям запросов, без независимых
человеческих оценщиков. Поэтому метрики характеризуют данный эксперимент
и не являются исчерпывающей оценкой всего корпуса.
